\section*{Part IV: Founder Decision Rights and Management Accountability}
\addcontentsline{toc}{section}{Part IV: Founder Decision Rights and Management Accountability}

\subsection{25. The delegation ladder}

Define levels: \textbf{inform} (employee decides within policy, then reports); \textbf{recommend} (employee develops alternatives, founder approves); \textbf{decide within limits} (budget and risk boundaries); \textbf{joint approval} (material or cross-functional commitments); and \textbf{board/stakeholder reserved} (statutory or constitutional matters). Assign a level to each recurring decision type.

Delegation fails when a person is accountable for outcomes without permission to make ordinary decisions. Founder control should focus on important boundaries and leadership appointments, not every invoice or lecture schedule.

\subsection{26. A decision-rights register}

At minimum specify who can approve hiring and termination, salary changes, discounts and refunds, contracts, customer claims, capital expenditure, loans, related-party dealings, IP licensing, data exports, research publication, admissions decisions, academic misconduct, property and new equity. State monetary thresholds, risk exceptions and escalation channels.

Review the register when the business changes, not merely at annual board renewal. Emergency exceptions should be logged and retrospectively reviewed without normalizing rule-breaking.

\subsection{27. Financial authority must separate custody and approval}

A founder-controlled startup can still maintain dual controls for substantial payments, bank account changes, contractor onboarding and payroll modifications. The person initiating a vendor payment should not unilaterally approve their own reimbursement. Reconcile bank statements and gateway receipts to invoices and revenue recognition. Use a documented cash forecast.

Segregation of duties can be achieved partly through external accounting review in tiny teams. The risk is not only fraud; it is error and false confidence in available cash.

\subsection{28. Board reporting: the minimum sufficient packet}

A monthly or quarterly packet should show liquidity and 13-week forecast, income statement and balance sheet, material liabilities, budget variance, cohort/product economics, receivables, quality/service incidents, material legal changes, headcount and concentrations, major decisions required, and conflicts disclosed. Include a short narrative of what management believed previously and what changed.

A founder who reports bad news early retains more credibility than 1 who provides polished numbers late. The board should make decisions, not merely hear an optimistic pitch.

\subsection{29. No approvals without recorded rationale}

For significant expenditures, document the problem, alternatives, expected benefits, financing, downside, owner, approval and review date. Record dissent and conditions. A 1-page decision memo is often sufficient for a young company. For related-party property, major VC terms or degree programs, the decision should receive independent legal and financial review.

An audit trail allows later leaders to understand why a decision was reasonable given information then - not only judge it with hindsight.

\subsection{30. Conflict of interest and related-party transactions}

A founder may own property leased to the company, hire relatives, invest in student startups, license personal IP, or transact with an entity they control. These may be legitimate, but only if disclosed, appropriately approved, fairly priced and compliant. Indian Companies Act provisions concerning director interests and related-party transactions are relevant, with exemptions and thresholds dependent on facts and amendments [G2].

Use independent market comparisons and recusal where required. ``The founder deserves the upside'' is not a substitute for a fair process when other shareholders, students, donors or the public have interests.
